\chapter{Stock Loan and Short Selling in Practice}\label{ch:m1:stock-loan-and-short-selling}

A fund is short a stock it believes is worth half its price, and pays 1\% a
year to borrow it. Others reach the same view; the shares that lenders offer
run out; within a week the fee is 80\%. The fund is now paying two tenths of
a percent of the position every day to keep a trade that may take a year to
work. Then a lender sells its shares and \omterm{def:m1:stock-loan-and-short-selling:recall}{recalls} the loan, the broker cannot
find replacements, and the fund must buy in a rising market along with every
other short. A long position can be wrong for years at no cost but patience.
A short position is rented, the rent is set daily by an auction the tenant
cannot see, and the landlord can end the lease overnight. \Cref{ch:m1:financing}
defined the loan; this chapter is about living with it.

\section{The lending chain}

\begin{definition}[Lendable supply and utilisation]\label{def:m1:stock-loan-and-short-selling:utilisation}
The \emph{lendable supply}\index{lendable supply} of a security is the
quantity that its owners have made available to lend through their agents.
\emph{Utilisation}\index{utilisation} is the fraction of it currently on
loan.
\end{definition}

The owners are pension funds, index funds, insurers and sovereign funds:
holders who will not sell next week and welcome a few basis points of
income. They lend through \emph{agent lenders}, usually their custodians, to
the \omterm{def:m1:financing:pb}{prime brokers} of \cref{def:m1:financing:pb}, who lend on to their
clients. \omterm{def:m1:financing:pb}{Prime brokers} also lend from their own clients' margin accounts
(\omterm{def:m1:financing:pb}{rehypothecation}) and borrow from one another. Long holders who have not
agreed to lend (most retail cash accounts, many active managers) are not in
the supply at all: for a typical large stock the \omterm{def:m1:stock-loan-and-short-selling:utilisation}{lendable supply} is a
fraction of the float, and what is actually on loan a small fraction of
that.

\begin{omfigure}
\begin{tikzpicture}[font=\small,
  box/.style={draw=omDef,rounded corners=2pt,minimum width=2.45cm,minimum height=1.05cm,align=center,fill=omDef!4},
  flow/.style={-{Stealth[length=2mm]},thick}]
\node[box] (own) at (0,0) {Beneficial\\owner};
\node[box] (agent) at (3.45,0) {Agent\\lender};
\node[box] (pb) at (6.9,0) {Prime\\broker};
\node[box] (hf) at (10.35,0) {Short\\seller};
\node[box,draw=omExo,fill=omExo!6] (buyer) at (10.35,-2.4) {Buyer in\\the market};
\foreach \a/\b in {own/agent,agent/pb,pb/hf}{
  \draw[flow,omDef] ([yshift=7pt]\a.east) -- ([yshift=7pt]\b.west);
  \draw[flow,omThm] ([yshift=-7pt]\b.west) -- ([yshift=-7pt]\a.east);}
\draw[flow,omDef] (hf.south) -- node[right,font=\footnotesize]{sells the shares} (buyer.north);
\node[font=\footnotesize,text=omDef] at (5.2,1.0) {shares (title passes; recallable at any time)};
\node[font=\footnotesize,text=omThm] at (5.2,-1.0) {collateral of about 102\%; the fee flows left, split at each step};
\end{tikzpicture}

\omcaption{The lending chain. The buyer at the end is a full owner, entitled
to vote, to receive dividends and to lend the shares again. The original
owner holds a claim on its borrower, receives a payment in lieu of each
dividend, and can \omterm{def:m1:stock-loan-and-short-selling:recall}{recall}.}\label{fig:m1:stock-loan-and-short-selling:chain}
\end{omfigure}

\begin{definition}[Short interest and days to cover]\label{def:m1:stock-loan-and-short-selling:si}
The \emph{short interest}\index{short interest} of a stock is the number of
its shares sold short and not yet covered, often quoted as a percentage of
the float. \emph{Days to cover}\index{days to cover} is short interest divided
by the average daily volume.
\end{definition}

Because the buyer of a shorted share can lend it again, the same share can
be shorted more than once, and \omterm{def:m1:stock-loan-and-short-selling:si}{short interest} can exceed 100\% of the float
without any failure to deliver: the longs simply hold, between them, the
float plus the \omterm{def:m1:stock-loan-and-short-selling:si}{short interest}.

\section{Fees, rebates and specials}

\begin{definition}[General collateral and special]\label{def:m1:stock-loan-and-short-selling:special}
A security is \emph{general collateral}\index{general collateral} when supply
so exceeds demand that it lends at a minimal fee, a few tenths of a percent a
year. It is \emph{special}\index{special} (or hard to borrow) when the fee is
materially higher; a special's \omterm{def:m1:financing:seclending}{rebate rate} is often negative: the borrower
receives no interest on its cash collateral and pays on top.
\end{definition}

The fee is the difference between the market interest rate and the rebate
(\cref{def:m1:financing:seclending}). It is set loan by loan, can be changed
by the lender on any day, and accrues on the \emph{current} market value of
the shares: when a short goes wrong the rent rises with the price. Fees are
negotiated bilaterally; no tape publishes them. Desks work from their \omterm{def:m1:financing:pb}{prime
brokers}' indications and from data vendors who pool the loan books of
lenders and brokers.

\begin{omfigure}
\begin{tikzpicture}
\begin{axis}[width=9.5cm,height=4.8cm,
  xlabel={utilisation (\%)},ylabel={borrow fee (\% a year)},
  tick label style={font=\small},label style={font=\small},
  xmin=0,xmax=100,ymin=0,ymax=85,grid=major,grid style={gray!25},table/col sep=comma]
\addplot[omDef,very thick] table[x=utilisation_pct,y=fee_pct]{figdata/markets-1/16-stock-loan-and-short-selling/fee.csv};
\node[font=\footnotesize,anchor=south] at (axis cs:38,4) {general collateral: 0.3\%};
\node[font=\footnotesize,anchor=east] at (axis cs:91,55) {special};
\end{axis}
\end{tikzpicture}

\omcaption{The chapter's illustrative fee curve: flat while supply is ample,
convex once three quarters of it is on loan. Real curves differ by stock and
by day; the shape, a long flat region and a steep end, is the robust
part.}\label{fig:m1:stock-loan-and-short-selling:fee}
\end{omfigure}

\begin{proposition}[The break-even holding period]\label{prop:m1:stock-loan-and-short-selling:breakeven}
A short seller expects the price to fall by a fraction $g$. With a borrow
fee $\phi$ a year, the fee consumes the whole expected gain after
\[
T^* \approx \frac{g}{\phi}\ \text{years},
\]
to first order (fee charged on the initial value, interest on the proceeds
ignored).
\end{proposition}

\begin{example}[Right, and unprofitable]\label{ex:m1:stock-loan-and-short-selling:right}
A thesis worth a 25\% fall has $T^* = 225$ days at a fee of 40\% and 112 days
at 80\%. \Cref{fig:m1:stock-loan-and-short-selling:carry} follows a short of
100\,000 shares at \$40 over a simulated year in which \omterm{def:m1:stock-loan-and-short-selling:utilisation}{utilisation} climbs from
60\% to 99\%. The stock first rises to \$53.52 on day 94, a loss of \$1.36
million. On day 190 it is at \$31.20 and the position is
\$789\,000 ahead. By day 250 the stock has recovered to \$38.36, still below
the entry, and the borrow has cost \$399\,000, three quarters of it in the
last sixty days: the trade has lost \$235\,000.
\end{example}

\begin{omfigure}
\begin{tikzpicture}
\begin{axis}[width=11cm,height=5.2cm,
  xlabel={day},ylabel={thousand dollars},
  tick label style={font=\small},label style={font=\small},
  xmin=1,xmax=250,ymin=-1500,ymax=1100,grid=major,grid style={gray!25},
  legend columns=3,legend style={font=\footnotesize,at={(0.5,-0.3)},anchor=north,draw=none,fill=none,/tikz/every even column/.append style={column sep=8pt}},
  table/col sep=comma]
\addplot[omExo,thick] table[x=day,y=price_pnl_k]{figdata/markets-1/16-stock-loan-and-short-selling/carry.csv};
\addplot[omThm,very thick,dashed] table[x=day,y=cost_k]{figdata/markets-1/16-stock-loan-and-short-selling/carry.csv};
\addplot[omDef,very thick] table[x=day,y=net_k]{figdata/markets-1/16-stock-loan-and-short-selling/carry.csv};
\legend{price P\&L,cumulative borrow cost,net}
\end{axis}
\end{tikzpicture}

\omcaption{A short of 100\,000 shares at \$40 while the stock goes \omterm{def:m1:stock-loan-and-short-selling:special}{special}.
The borrow cost is negligible for half a year and then takes everything.
Data: the tutorial's simulation.}\label{fig:m1:stock-loan-and-short-selling:carry}
\end{omfigure}

\section{Recalls, buy-ins and the rules}

\begin{definition}[Recall and buy-in]\label{def:m1:stock-loan-and-short-selling:recall}
A \emph{recall}\index{recall} is the lender's demand for the return of lent
securities, typically because it has sold them or wants to vote. If the
borrower's broker cannot replace the loan from another lender and the
borrower does not return the shares in time, the lender (or the broker)
executes a \emph{buy-in}\index{buy-in}: it buys the shares in the market at
the borrower's expense.
\end{definition}

\begin{definition}[Naked short sale]\label{def:m1:stock-loan-and-short-selling:naked}
A \emph{naked short sale}\index{naked short sale} is a \omterm{def:m1:financing:short}{short sale} made
without having borrowed, or arranged to borrow, the security in time for
settlement, so that the seller fails to deliver
(\cref{ch:m1:clearing-and-settlement}).
\end{definition}

\begin{dated}{2026-09}{The rules on each side of the Atlantic}\label{dat:m1:stock-loan-and-short-selling:rules}
\textbf{United States, Regulation SHO.} \emph{Locate} (Rule 203(b)(1)): before
effecting a \omterm{def:m1:financing:short}{short sale} a \omterm{def:m1:the-sell-side:sell-side}{broker-dealer} must have reasonable grounds to
believe the security can be borrowed and delivered on the due date.
\emph{Close-out} (Rule 204): a fail from a \omterm{def:m1:financing:short}{short sale} must be closed out by
the beginning of regular trading hours on the settlement day following the
settlement date; fails from long sales and from bona fide market making have
until the third settlement day following it. \emph{Threshold securities}:
fails of at least 10\,000 shares and 0.5\% of the shares outstanding for five
consecutive settlement days; fails persisting for 13 consecutive settlement
days must be closed out at once. \emph{Price test} (Rule 201): after a
fall of 10\% in a day, \omterm{def:m1:financing:short}{short sales} are restricted to prices above the national
best bid for the rest of that day and the next. The regulator's own summary
notes that naked short selling is not in itself a violation.

\textbf{European Union, Short Selling Regulation.} Net short positions in
shares are notified to the national authority from 0.1\% of the issued share
capital and at each further 0.1\%, and published from 0.5\%. Uncovered \omterm{def:m1:financing:short}{short
sales} of shares and of sovereign debt are prohibited, and national
authorities may impose temporary bans in a crisis.
\end{dated}

Both regimes make an exception for market making: in the United States a
\omterm{def:m1:the-sell-side:sell-side}{broker-dealer} engaged in bona fide market making need not locate before
selling, since a firm that must quote both sides continuously cannot consult
a stock-loan desk before each sale; the exception does not cover speculative
selling or one-sided quoting. The exception, and the longer close-out, are
part of the economics of the \omterm{def:m1:what-a-trading-firm-does:market-maker}{market makers} of
\cref{ch:m1:what-a-trading-firm-does}.

\section{Squeezes}

\begin{definition}[Short squeeze]\label{def:m1:stock-loan-and-short-selling:squeeze}
A \emph{short squeeze}\index{short squeeze} is a rise in price amplified by
the purchases of short sellers who cover because of losses, \omterm{def:m1:financing:margincall}{margin calls},
\omterm{def:m1:stock-loan-and-short-selling:recall}{recalls} or risk limits, each purchase pushing the price further against
those who remain.
\end{definition}

\begin{proposition}[The covering cascade]\label{prop:m1:stock-loan-and-short-selling:cascade}
Let the shorts' covering thresholds, expressed as returns from today's price,
be spread uniformly between $a$ and $b$, and let $k$ be the price impact, as a
return, of all of them covering. After an exogenous shock $s > a$ the price
settles at the return
\[
x \;=\; \frac{s - ka/(b-a)}{1 - k/(b-a)} \qquad\text{if } k < b-a \text{ and } x < b,
\]
and at $x = s + k$, with every short covered, otherwise. The multiplier
$1/(1-k/(b-a))$ is unbounded as $k \to b-a$.
\end{proposition}
\begin{proof}
The fraction covered at return $x$ is $F(x) = (x-a)/(b-a)$ on $[a,b]$, and the
price satisfies $x = s + kF(x)$. Solve the linear equation. If $k \ge b-a$
each cover triggers at least one more and the iteration $x \leftarrow s +
kF(x)$ stops only at $F = 1$.
\end{proof}

\begin{omfigure}
\begin{tikzpicture}
\begin{axis}[width=10.5cm,height=5.4cm,
  xlabel={initial shock (\%)},ylabel={final rise (\%)},
  tick label style={font=\small},label style={font=\small},
  xmin=0,xmax=40,ymin=0,ymax=105,grid=major,grid style={gray!25},
  legend columns=4,legend style={font=\footnotesize,at={(0.5,-0.3)},anchor=north,draw=none,fill=none,/tikz/every even column/.append style={column sep=6pt}},
  table/col sep=comma]
\addplot[omExo,thick,dotted] coordinates {(0,0) (40,40)};
\addplot[omDef,very thick] table[x=shock_pct,y=k20]{figdata/markets-1/16-stock-loan-and-short-selling/squeeze.csv};
\addplot[omThm,very thick,dashed] table[x=shock_pct,y=k40]{figdata/markets-1/16-stock-loan-and-short-selling/squeeze.csv};
\addplot[black,very thick,dashdotted] table[x=shock_pct,y=k60]{figdata/markets-1/16-stock-loan-and-short-selling/squeeze.csv};
\legend{no shorts,$k = 20\%$,$k = 40\%$,$k = 60\%$}
\end{axis}
\end{tikzpicture}

\omcaption{The covering cascade with thresholds between $+10\%$ and $+60\%$.
With $k = 20\%$ a shock is amplified by $5/3$; with $k = 40\%$ by 5, until at a
shock of 20\% everyone has covered; with $k = 60\% > b - a$ the first cover
sets off all the others. Data:
\cref{prop:m1:stock-loan-and-short-selling:cascade}.}\label{fig:m1:stock-loan-and-short-selling:squeeze}
\end{omfigure}

The impact $k$ grows with \omterm{def:m1:stock-loan-and-short-selling:si}{short interest} relative to the shares that can
actually be bought: \omterm{def:m1:stock-loan-and-short-selling:si}{days to cover}, and \omterm{def:m1:stock-loan-and-short-selling:si}{short interest} over the \emph{free}
float, are the two numbers to watch. Two episodes mark the extremes.

\begin{example}[Wolfsburg, October 2008]\label{ex:m1:stock-loan-and-short-selling:vw}
On Sunday 26 October 2008 a carmaker announced that it held 42.6\% of another
carmaker's ordinary shares and cash-settled options on a further 31.5\%. A
German state held 20\%. About 13\% of the shares were sold short, and 5.9\%
remained in the float. The price, \euro211 on the Friday, passed \euro1\,005
on Tuesday 28 October, briefly making the company the most valuable listed
company in the world. The options were the point: being cash-settled they
fell outside the disclosure rules of the time, and whatever shares their
writers held as hedges were not for sale.
\end{example}

\begin{example}[January 2021]\label{ex:m1:stock-loan-and-short-selling:gme}
\omterm{def:m1:stock-loan-and-short-selling:si}{Short interest} in a US video-game retailer reached 122.97\% of its float in
January 2021. From an intraday low on 8 January to an intraday high of \$483
on 28 January the price rose about 2\,700\%, then fell more than 86\% by the
end of the first week of February. The regulator's staff report found that
short sellers' covering coincided with parts of the rise but was a small
fraction of total buying, and concluded that positive sentiment, not
buying-to-cover, sustained the rise: a squeeze was part of the story and not
the whole of it.
\end{example}

\section{Dividends and tax}

The borrower of a share owes the lender a \emph{manufactured dividend} for
every dividend paid during the loan. The lender is economically whole; but it
is no longer the legal owner on the \omterm{def:m1:shares-corporate-actions-indices:dividend}{record date}, and where a withholding tax
applies, the tax position of owner, borrower and buyer can differ. Lending
shares over the dividend date to a holder with a better tax position, and
sharing the saving, is an old and in many jurisdictions contested trade.

Its criminal extreme was \emph{cum-ex}: \omterm{def:m1:financing:short}{short sales} arranged around the
dividend date so that the withholding tax paid once appeared to have been
paid to two owners, both of whom reclaimed it. On 28 July 2021 Germany's
Federal Court of Justice confirmed that claiming refunds of tax that had
never been paid is tax evasion, a criminal offence. For a desk the lesson is
narrow and absolute: a dividend trade whose profit is a tax refund is a legal
question before it is a trading question.

\section{Tutorial: the cost of a short as the stock goes special}\label{tut:m1:stock-loan-and-short-selling:cost}

\begin{tutorial}
\textbf{Goal.} Carry a short through a year in which the borrow tightens, and
reproduce the covering cascade. \textbf{End state:} the three data figures of
this chapter.
\begin{tutsteps}
\item \textbf{A fee curve.} Flat, then convex.
\omcode{code/markets-1/16-stock-loan-and-short-selling/python/short_cost.py}{5}{10}{An illustrative borrow fee as a function of utilisation.}\label{lst:m1:stock-loan-and-short-selling:fee}
\item \textbf{Carry.} The fee accrues on the day's market value, actual/360.
\omcode{code/markets-1/16-stock-loan-and-short-selling/python/short_cost.py}{13}{17}{Price P\&L and cumulative borrow cost of a short.}\label{lst:m1:stock-loan-and-short-selling:carry}
\item \textbf{The cascade.} A closed form and an iteration that must agree.
\omcode{code/markets-1/16-stock-loan-and-short-selling/python/short_cost.py}{29}{37}{The equilibrium of the covering cascade.}\label{lst:m1:stock-loan-and-short-selling:squeeze}
\end{tutsteps}
\textbf{What to change next.} Make \omterm{def:m1:stock-loan-and-short-selling:utilisation}{utilisation} respond to the price (shorts
add when the stock rises on no news, lenders sell into strength and \omterm{def:m1:stock-loan-and-short-selling:recall}{recall}).
Replace the uniform thresholds by a distribution concentrated near round
losses and look at what happens to the multiplier.
\end{tutorial}

\section{Build: the borrow book}\label{bld:m1:stock-loan-and-short-selling:borrow}

\begin{build}
\bfield{Purpose} The financing calculator of \cref{ch:m1:financing} takes
borrow fees as an input. This component produces them, answers locate
requests, and tells the firm what it must buy in when a lender \omterm{def:m1:stock-loan-and-short-selling:recall}{recalls}.
\bfield{Interface} \texttt{FeeCurve(gc, kink, max\_fee).fee(utilisation)};
\texttt{BorrowBook.locate(sym, qty)}, \texttt{borrow}, \texttt{give\_back},
\texttt{recall(sym, qty)} returning the firm's \omterm{def:m1:stock-loan-and-short-selling:recall}{buy-in} quantity,
\texttt{utilisation(sym)} and \texttt{borrow\_fees()}, a dictionary accepted
as is by \texttt{firm\_financing.accrue}.
\bfield{Rules} The firm's own borrow counts in \omterm{def:m1:stock-loan-and-short-selling:utilisation}{utilisation}: a large short
moves its own fee. A \omterm{def:m1:stock-loan-and-short-selling:recall}{recall} that leaves loans above supply is shared pro rata
between the firm and the rest of the market, rounded against the firm. A
locate for an unknown symbol returns zero, never an error: no locate, no
short.
\bfield{Acceptance tests} \texttt{code/firm/borrow/tests/}: locates and fee
levels; the firm's borrow moving the fee; a \omterm{def:m1:stock-loan-and-short-selling:recall}{recall} with a pro-rata \omterm{def:m1:stock-loan-and-short-selling:recall}{buy-in};
the hand-off to the financing accrual.
\bfield{Stretch} Term loans at a fixed fee against overnight loans at a
floating one: the choice every stock-loan desk faces before an event.
\end{build}

\begin{omsources}
\item US Securities and Exchange Commission, \emph{Key Points About
Regulation SHO}, investor publication.
\item US Securities and Exchange Commission, \emph{Staff Report on Equity and
Options Market Structure Conditions in Early 2021}, 14 October 2021.
\item European Securities and Markets Authority, \emph{Short selling}, policy
page; Regulation (EU) No 236/2012.
\item F.~Allen, M.~Haas, E.~Nowak and A.~Tengulov, ``Market efficiency and
limits to arbitrage: evidence from the Volkswagen short squeeze'',
\emph{Journal of Financial Economics} 142 (2021); ``The case of Volkswagen'',
\emph{The Hedge Fund Journal}.
\item Bundesgerichtshof, press release 146/2021 on judgment 1 StR 519/20 of
28 July 2021.
\item D.~Duffie, N.~G\^arleanu and L.~H.~Pedersen, ``Securities lending,
shorting, and pricing'', \emph{Journal of Financial Economics} 66 (2002).
\end{omsources}

\section{Exercises}

\begin{exercise}[$\star$]\label{exo:m1:stock-loan-and-short-selling:1}
The overnight rate is 4.3\% and a \omterm{def:m1:financing:pb}{prime broker} quotes a rebate of $-11.7\%$ on
a stock. Give the borrow fee and the daily cost of a \$3 million short.
\end{exercise}

\begin{exercise}[$\star$]\label{exo:m1:stock-loan-and-short-selling:2}
A stock has 12 million shares of \omterm{def:m1:stock-loan-and-short-selling:utilisation}{lendable supply}, 10.2 million of them on
loan. Give the \omterm{def:m1:stock-loan-and-short-selling:utilisation}{utilisation} and the fee on the chapter's curve.
\end{exercise}

\begin{exercise}[$\star$]\label{exo:m1:stock-loan-and-short-selling:3}
A stock has a float of 60 million shares, a \omterm{def:m1:stock-loan-and-short-selling:si}{short interest} of 18 million and
an average volume of 2.4 million a day. Give the \omterm{def:m1:stock-loan-and-short-selling:si}{short interest} as a
percentage of float and the \omterm{def:m1:stock-loan-and-short-selling:si}{days to cover}. Which of the two matters more to
a short seller worried about a squeeze, and why?
\end{exercise}

\begin{exercise}[$\star\star$]\label{exo:m1:stock-loan-and-short-selling:4}
An analyst expects a 25\% fall. Give the break-even holding period at fees
of 40\% and 80\%. She expects the fall to happen ``some time in the next two
years'' with uniform probability. Roughly what fraction of her scenarios make
money at 40\%?
\end{exercise}

\begin{exercise}[$\star\star$]\label{exo:m1:stock-loan-and-short-selling:5}
A company has a float of 100 shares, all held by A, who lends them. Construct
a sequence of loans and sales after which \omterm{def:m1:stock-loan-and-short-selling:si}{short interest} is 123\% of the
float and nobody has failed to deliver. How many shares do the longs hold in
total? Who votes?
\end{exercise}

\begin{exercise}[$\star\star$]\label{exo:m1:stock-loan-and-short-selling:6}
A \omterm{def:m1:financing:short}{short sale} executed on Monday 14 September 2026 fails to settle. Using
\cref{dat:m1:stock-loan-and-short-selling:rules} and one-day settlement,
give the deadline for the close-out. What would it be for a long sale that
failed?
\end{exercise}

\begin{exercise}[$\star\star\star$]\label{exo:m1:stock-loan-and-short-selling:7}
\emph{Coding.} From \texttt{carry.csv} (or by running the figure script)
report the net P\&L at its maximum with its day, the final price P\&L, the
final cumulative borrow cost, and the share of that cost incurred after day
190.
\end{exercise}

\begin{exercise}[$\star\star\star$]\label{exo:m1:stock-loan-and-short-selling:8}
\emph{Find the flaw.} ``\omterm{def:m1:stock-loan-and-short-selling:si}{Short interest} is 140\% of the float. More shares have
been sold than exist. This proves illegal naked shorting with counterfeit
shares.'' Explain why the inference fails, and say what data \emph{would} be
evidence of naked shorting.
\end{exercise}

\section{Problem: The Squeeze}

\begin{problem}[{Weekend problem --- a crowded short meets good news}]\label{pb:m1:stock-loan-and-short-selling:1}
A fund is short 2 million shares of Z at \$20. Z has a float of 50 million
shares, a \omterm{def:m1:stock-loan-and-short-selling:si}{short interest} of 20 million and an average volume of 4 million a
day. The fund expects a fall of 30\%. The borrow fee is 12\%.

\medskip\noindent\textbf{Part I --- The carry.}
\begin{enumerate}
\item Give the \omterm{def:m1:stock-loan-and-short-selling:si}{short interest} as a percentage of float, the \omterm{def:m1:stock-loan-and-short-selling:si}{days to cover},
and the fund's share of the \omterm{def:m1:stock-loan-and-short-selling:si}{short interest}.
\item Give the daily borrow cost.
\item Give the break-even holding period.
\item The fee reprices to 60\%. Give the new daily cost and break-even.
\item The fund's average holding period for such trades is nine months. What
does question 4 imply?
\end{enumerate}

\medskip\noindent\textbf{Part II --- The news.}
Z rises 25\% on a takeover rumour. The shorts' covering thresholds are
uniform between $+10\%$ and $+60\%$.
\begin{enumerate}[resume]
\item With $k = 20\%$, where does the price settle, and what fraction of the
shorts has covered?
\item With $k = 40\%$?
\item For $k = 40\%$, find the smallest shock after which every short covers.
\item Give the fund's \omterm{def:m1:pnl-and-positions:mtm}{mark-to-market} loss at the price of question 7.
\item Its \omterm{def:m1:financing:pb}{prime broker} raises the margin on the short from 30\% to 100\% of
market value. Give the additional margin at that price.
\end{enumerate}

\medskip\noindent\textbf{Part III --- The \omterm{def:m1:stock-loan-and-short-selling:recall}{recall}.}
\begin{enumerate}[resume]
\item A lender \omterm{def:m1:stock-loan-and-short-selling:recall}{recalls}; the broker cannot replace 500\,000 of the fund's
shares and buys them in at \$33. Give the realised loss.
\item Why do \omterm{def:m1:stock-loan-and-short-selling:recall}{recalls} cluster exactly when the price has risen?
\item The fund had the choice, a month earlier, of a three-month term borrow
at a fixed 20\%. Give its cost for three months and the argument for paying
it.
\item Propose a position in listed options with the same view as the short
and no borrow. What does it cost instead, and where has the borrow fee gone?
\end{enumerate}

\medskip\noindent\textbf{Part IV --- Judgement.}
\begin{enumerate}[resume]
\item Which of the numbers of question 1 would have warned the fund?
\item A colleague says the squeeze ``proves manipulation''. Using the staff
report of \cref{ex:m1:stock-loan-and-short-selling:gme}, explain why that
needs evidence of its own.
\item The rumour proves false and Z falls to \$13 within a quarter. Was the
fund right?
\item How should short positions be sized, given Parts II and III?
\item State the \emph{named result}: the break-even holding period of the
short after the fee repriced.
\item In one sentence: what does a short seller rent, and from whom?
\end{enumerate}
\end{problem}

\section{Interview questions}

\begin{interviewq}[$\star$ \iqroles{trader, researcher}]\label{iq:m1:stock-loan-and-short-selling:1}
Walk me through what happens, operationally, when a fund shorts a stock.
\end{interviewq}

\begin{interviewq}[$\star$ \iqroles{trader, researcher}]\label{iq:m1:stock-loan-and-short-selling:2}
How can \omterm{def:m1:stock-loan-and-short-selling:si}{short interest} exceed 100\% of the float?
\end{interviewq}

\begin{interviewq}[$\star\star$ \iqroles{researcher, mle}]\label{iq:m1:stock-loan-and-short-selling:3}
Your backtest shorts the bottom decile of a signal and earns 9\% a year on
the short leg. What is missing?
\end{interviewq}

\begin{interviewq}[$\star\star$ \iqroles{trader, researcher}]\label{iq:m1:stock-loan-and-short-selling:4}
A stock's borrow fee has gone from 2\% to 45\% in a week. What does that tell
you, and what does it not?
\end{interviewq}

\begin{interviewq}[$\star\star$ \iqroles{trader, bank}]\label{iq:m1:stock-loan-and-short-selling:5}
Why are \omterm{def:m1:what-a-trading-firm-does:market-maker}{market makers} exempt from the locate requirement, and what stops
them from abusing it?
\end{interviewq}

\begin{interviewq}[$\star\star\star$ \iqroles{researcher, trader}]\label{iq:m1:stock-loan-and-short-selling:6}
Put options on a hard-to-borrow stock look expensive against calls: put-call
parity seems violated. Is there an arbitrage?
\end{interviewq}
